\documentclass[10pt,a4paper]{article}

\usepackage[utf8]{inputenc}
\usepackage[T1]{fontenc}
\usepackage[spanish,es-nodecimaldot]{babel}
\usepackage[margin=2.2cm]{geometry}
\usepackage{amsmath,amssymb}
\usepackage{booktabs,array}
\usepackage{tikz}
\usepackage{microtype}
\usepackage[hidelinks]{hyperref}
\setlength{\parindent}{0pt}
\setlength{\parskip}{4pt}
\renewcommand{\arraystretch}{1.08}
\newcommand{\ruta}[1]{\ensuremath{#1}}

\begin{document}

\begin{center}
  {\Large\bfseries Práctica VRP}\\[4pt]
  Logística y Servicios\\[2pt]
  Modelo capacitado y heurística de Clarke y Wright
\end{center}

\section{Datos y modelo matemático}

La instancia contiene un depósito, situado en el nodo 0, y cinco clientes.
Los datos proceden de la hoja \texttt{VRP 5} del fichero
\texttt{Tema 4 VRP\_datos.xlsx}. La capacidad del vehículo es $Q=25$.

\begin{center}
\begin{tabular}{rrrr}
\toprule
Nodo & Coordenada $X$ & Coordenada $Y$ & Demanda $d_i$\\
\midrule
0 & 20 & 4 & 0\\
1 & $-7$ & 9 & 8\\
2 & $-3$ & 10 & 6\\
3 & $-17$ & $-9$ & 8\\
4 & 8 & $-16$ & 7\\
5 & 7 & $-20$ & 3\\
\bottomrule
\end{tabular}
\end{center}

El coste del arco $(i,j)$ es la distancia euclídea
\[
 c_{ij}=\sqrt{(X_i-X_j)^2+(Y_i-Y_j)^2}.
\]
No se redondean las distancias durante los cálculos. Los costes se expresan
en unidades de distancia, puesto que no se especifican una escala geográfica
ni una tarifa de transporte. La demanda total es $32$, por lo que se necesitan
al menos $\lceil32/25\rceil=2$ viajes. Un mismo vehículo puede realizarlos
sucesivamente y recargar al regresar al depósito.

Sean $N=\{0,1,2,3,4,5\}$ y $V=N\setminus\{0\}$. La variable binaria
$x_{ij}$ vale 1 si se recorre el arco $(i,j)$. Para cada cliente, $u_i$
es una variable auxiliar de carga acumulada que permite imponer la capacidad
y eliminar subciclos; sus restricciones pueden admitir holgura.

\begin{align}
 \min\quad & z=\sum_{i\in N}\sum_{\substack{j\in N\\j\ne i}}c_{ij}x_{ij}
 \label{eq:obj}\\
 \text{s.a.}\quad
 & \sum_{\substack{j\in N\\j\ne i}}x_{ij}=1 && \forall i\in V,\\
 & \sum_{\substack{i\in N\\i\ne j}}x_{ij}=1 && \forall j\in V,\\
 & \sum_{j\in V}x_{0j}=\sum_{i\in V}x_{i0},\\
 & u_i-u_j+Qx_{ij}\le Q-d_j && \forall i,j\in V,\ i\ne j,\label{eq:mtz}\\
 & d_i\le u_i\le Q && \forall i\in V,\\
 & x_{ij}\in\{0,1\} && \forall i,j\in N,\ i\ne j.
\end{align}

Las restricciones de entrada y salida garantizan una visita por cliente.
El balance del depósito iguala las salidas y los regresos, sin fijar el número
de viajes. Si $x_{ij}=1$, la restricción~\eqref{eq:mtz} implica
$u_j\ge u_i+d_j$. Por tanto, las restricciones MTZ impiden subciclos entre
clientes y, junto con las cotas de $u_i$, limitan la carga de cada ruta.
En la matriz de decisión del Excel también se impone $x_{ii}=0$.

\newpage
\section{Solución óptima}

El modelo se ha resuelto con el complemento Solver de Microsoft Excel
Online. Solver ha indicado que se satisfacen todas las restricciones y las
condiciones de optimalidad. La matriz de decisión obtenida determina los
siguientes viajes:

\begin{center}
\begin{tabular}{lrr}
\toprule
Ruta & Carga & Coste\\
\midrule
\ruta{0-3-1-2-0} & 22 & $87{,}701438$\\
\ruta{0-4-5-0}   & 10 & $54{,}741601$\\
\midrule
Total & 32 & $142{,}443039$\\
\bottomrule
\end{tabular}
\end{center}

Las cargas son $d_3+d_1+d_2=8+8+6=22$ y $d_4+d_5=7+3=10$.
Ambas rutas respetan la capacidad y atienden a todos los clientes una sola vez.
Los arcos seleccionados son $(0,3)$, $(3,1)$, $(1,2)$, $(2,0)$,
$(0,4)$, $(4,5)$ y $(5,0)$.

Solver ha asignado $u_3=8$, $u_1=19$, $u_2=25$, $u_4=7$ y $u_5=25$.
Estos valores satisfacen las restricciones MTZ, aunque algunos contienen
holgura. No son las demandas efectivamente entregadas: las cargas acumuladas
reales son $8,16,22$ en el primer viaje y $7,10$ en el segundo.

\begin{figure}[h]
\centering
\begin{tikzpicture}[x=0.22cm,y=0.22cm,
  punto/.style={circle,draw=black,fill=white,inner sep=2.2pt},
  deposito/.style={rectangle,draw=black,fill=black,inner sep=2.5pt}]
  \coordinate (n0) at (20,4);
  \coordinate (n1) at (-7,9);
  \coordinate (n2) at (-3,10);
  \coordinate (n3) at (-17,-9);
  \coordinate (n4) at (8,-16);
  \coordinate (n5) at (7,-20);
  \draw[thick] (n0)--(n3)--(n1)--(n2)--(n0);
  \draw[thick,dashed] (n0)--(n4)--(n5)--(n0);
  \node[deposito,label=right:0] at (n0) {};
  \node[punto,label=above:1] at (n1) {};
  \node[punto,label=above:2] at (n2) {};
  \node[punto,label=left:3] at (n3) {};
  \node[punto,label=right:4] at (n4) {};
  \node[punto,label=below:5] at (n5) {};
\end{tikzpicture}
\caption{Rutas óptimas sobre las coordenadas de la instancia. Trazo continuo:
$0-3-1-2-0$. Trazo discontinuo: $0-4-5-0$.}
\end{figure}

\subsection{Verificación exacta del valor óptimo}

Como comprobación independiente de Solver, el valor óptimo se ha verificado
mediante enumeración completa. Se examinan
las $5!=120$ permutaciones de los clientes y, para cada una, las $2^4=16$
formas de introducir cortes entre clientes consecutivos. Cada bloque
resultante se interpreta como un viaje que sale del depósito y regresa a él.

Este procedimiento cubre todos los órdenes de visita y todas las particiones
en rutas. Entre las $120\cdot16=1.920$ representaciones, incluidas las
equivalentes por simetría o por orden de los viajes, $1.704$ son factibles.
El menor coste entre ellas es
\[
  \boxed{z^*=142{,}44303899076746}.
\]

El coste obtenido con Solver coincide con este mínimo. La primera ruta es
la inversa de $0-2-1-3-0$, con el mismo coste por la simetría de las distancias.
El orden de ejecución de los dos viajes tampoco modifica el objetivo.

\newpage
\section{Heurística de Clarke y Wright}

La solución inicial consiste en cinco viajes individuales $0-i-0$.
El ahorro al unir los clientes $i$ y $j$ es
\[
 s_{ij}=c_{0i}+c_{0j}-c_{ij}.
\]
Los pares se consideran por ahorro decreciente. Solo se unen extremos de
rutas y la carga resultante no puede superar $Q$. En caso de empate se usa
el orden ascendente de $i$ y después de $j$. En esta instancia no hay empates.

\begin{center}
\begin{tabular}{crr@{\qquad}crr}
\toprule
Orden & Par & Ahorro & Orden & Par & Ahorro\\
\midrule
1 & $1-2$ & $47{,}105683$ & 6  & $3-4$ & $36{,}579641$\\
2 & $4-5$ & $46{,}495390$ & 7  & $1-5$ & $22{,}551264$\\
3 & $1-3$ & $46{,}085143$ & 8  & $1-4$ & $21{,}628109$\\
4 & $3-5$ & $40{,}111274$ & 9  & $2-5$ & $19{,}441640$\\
5 & $2-3$ & $39{,}386224$ & 10 & $2-4$ & $18{,}862348$\\
\bottomrule
\end{tabular}
\end{center}

\subsection{Esquema secuencial}

Se completa una ruta antes de comenzar otra. El par $1-2$ inicia la primera
ruta, con carga 14. El par $4-5$ no toca la ruta activa y se omite de momento.
El par $1-3$ permite añadir el cliente 3 por el extremo 1, con carga 22.
Después, $3-5$ incorpora el cliente 5 y alcanza la capacidad 25. Tras cada
unión se vuelve a recorrer la lista para buscar una extensión factible.
El cliente 4 queda fuera y requiere un viaje individual.

\begin{center}
\begin{tabular}{lrr}
\toprule
Ruta & Carga & Coste\\
\midrule
\ruta{0-5-3-1-2-0} & 25 & $102{,}179540$\\
\ruta{0-4-0} & 7 & $46{,}647615$\\
\midrule
Total & 32 & $148{,}827155$\\
\bottomrule
\end{tabular}
\end{center}

\subsection{Esquema paralelo}

Las rutas se construyen simultáneamente. Se aceptan, en este orden,
los pares $1-2$, $4-5$ y $1-3$. Quedan las rutas $0-2-1-3-0$ y $0-4-5-0$,
con cargas 22 y 10, respectivamente. Unirlas exigiría cargar 32 unidades y
se rechaza. El par $2-3$ ya pertenece a la misma ruta; los demás pares
restantes exceden la capacidad o afectan a nodos interiores.
El coste total del esquema paralelo es $142{,}443039$.

\section{Comparación de resultados}

La desviación porcentual relativa respecto al óptimo se calcula como
\[
 \operatorname{RPD}(\%)=100\,\frac{z_H-z^*}{z^*}.
\]
\begin{center}
\begin{tabular}{lrrr}
\toprule
Método & Coste & RPD & Viajes\\
\midrule
Modelo con Excel Solver & $142{,}443039$ & 0\,\% & 2\\
Clarke y Wright secuencial & $148{,}827155$ & 4,4819\,\% & 2\\
Clarke y Wright paralelo & $142{,}443039$ & 0\,\% & 2\\
\bottomrule
\end{tabular}
\end{center}

La variante paralela alcanza el óptimo de esta instancia. La secuencial
cuesta $6{,}384116$ unidades más, porque asigna el cliente 5 a la primera ruta
y deja al cliente 4 solo. En el esquema paralelo se forma antes el grupo
$4-5$, lo que permite obtener las dos rutas óptimas. Este resultado no implica
que la variante paralela sea siempre óptima ni superior a la secuencial.

\smallskip
{\footnotesize\textbf{Fuentes:} \textit{Práctica VRP.pdf}, p.\ 1;
\texttt{Tema 4 VRP\_datos.xlsx}, hoja \texttt{VRP 5}; y
\textit{Tema4-RoutingProblemsVRP\_Clase 5.pdf}, diapositivas 7--9 y 13--20.}

\end{document}
